\documentclass[10pt, a4paper]{article}

% ==========================================
% PRÉAMBULE
% ==========================================
\usepackage[utf8]{inputenc}
\usepackage[T1]{fontenc}
\usepackage[french]{babel}
\usepackage{geometry}
\usepackage{hyperref}
\usepackage{fancyhdr}
\usepackage{xcolor}
\usepackage{booktabs}
\usepackage{array}
\usepackage{longtable}
\usepackage{enumitem}
\usepackage{titlesec}

\geometry{top=2cm, bottom=2cm, left=2cm, right=2cm}
\setlist{nosep, leftmargin=1.2em}

\hypersetup{
    colorlinks=true,
    linkcolor=black,
    filecolor=magenta,
    urlcolor=blue,
    pdftitle={Cahier des Charges CARIBOOKS - Alexandre REY},
}

\pagestyle{fancy}
\fancyhf{}
\fancyhead[L]{Cahier des Charges — CARIBOOKS}
\fancyhead[R]{Alexandre REY - ESTIAM}
\fancyfoot[C]{\thepage}

\titlespacing*{\section}{0pt}{10pt}{4pt}
\titlespacing*{\subsection}{0pt}{6pt}{2pt}

\newcolumntype{L}[1]{>{\raggedright\arraybackslash}p{#1}}

\title{\textbf{Cahier des Charges — CARIBOOKS}\\[4pt]
\large\textit{Plateforme e-commerce solidaire de vente de livres d'occasion — Recyclerie Caritas}}
\author{Alexandre REY \\ E3 DAD — Concepteur Développeur d'Applications (CDA) \\ ESTIAM — Encadrant : Mhand Boufala}
\date{Année scolaire 2025-2026}

\begin{document}
\maketitle
\thispagestyle{fancy}

\noindent{\small Ce document formalise le besoin exprimé par le commanditaire (Recyclerie
Caritas) et les spécifications attendues de la plateforme CARIBOOKS. Il ne traite pas des choix
de solution technique, documentés dans le mémoire de projet (\texttt{main.tex}).}

\section{Contexte et objectifs}

L'association \textbf{Caritas} facilite la réinsertion professionnelle en Suisse romande via
la vente d'objets de seconde main. Une partie du stock de livres d'occasion ne trouve pas
preneur en magasin. La commanditaire (bénévole référente, Product Owner informelle) souhaite
digitaliser cette activité avec un outil simple pour des bénévoles non technophiles.

\noindent\textbf{Enjeux :} économique (limiter les invendus), social (soutenir la mission de
Caritas), opérationnel (saisie rapide sans compétence technique), financier (pas d'abonnement
SaaS récurrent).

\noindent\textbf{Public cible :} associations caritatives, étudiants et grand public résidant
en Suisse, achats en CHF avec livraison nationale.

\noindent\textbf{Objectifs opérationnels :} (1) ajout d'un livre en moins de 2 minutes via
scan ISBN ; (2) recherche par filtres combinés ; (3) paiement CHF sécurisé, livraison
restreinte à la Suisse ; (4) back-office autonome pour les bénévoles.

\noindent\textbf{Critères de réussite :} mise en vente < 2 min ; paiement CHF sans erreur
bloquante en sandbox PostFinance ; zéro prix hors CHF ou livraison hors Suisse ; tests
automatisés sur les parcours critiques.

\section{Périmètre du projet}

\begin{longtable}{@{}L{1.2cm}L{3.5cm}L{9.8cm}@{}}
\toprule
\textbf{Réf.} & \textbf{Fonctionnalité} & \textbf{Description} \\
\midrule
\endhead
F01 & Inscription / connexion & Compte, authentification, rôles utilisateur / administrateur. \\
F02 & Ajout d'un livre & Scan ISBN (préremplissage) ou saisie manuelle en secours. \\
F03 & Prix et état & Prix en CHF, état parmi \textit{comme neuf}, \textit{bon état}, \textit{état correct}. \\
F04 & Gestion des annonces & Publication, modification, suppression. \\
F05 & Catalogue et recherche & Filtres combinés (titre, auteur, genre, prix, état). \\
F06 & Fiche détaillée & Informations complètes du livre et du vendeur. \\
F07 & Panier et réservation de stock & Réservation temporaire, libération automatique en cas d'abandon. \\
F08 & Paiement sécurisé & CHF via PostFinance (TWINT, Visa, Mastercard). \\
F09 & Livraison Suisse & Envoi postal (tarif fixe) ou retrait en magasin (Click \& Collect). \\
F10 & Back-office bénévole & Gestion des annonces et des commandes. \\
F11 & Conformité RGPD / nLPD & Export et effacement des données personnelles. \\
\bottomrule
\caption{Périmètre fonctionnel du MVP}
\end{longtable}

\noindent\textbf{Hors périmètre :} messagerie acheteur/vendeur, suivi logistique détaillé,
estimation automatique du prix de revente, recommandations et avis.

\noindent\textbf{Frontières :} intègre OpenLibrary (métadonnées), PostFinance (paiement) et
Google OAuth (connexion optionnelle) ; ne gère ni comptabilité ni logistique physique
(aucune API de transporteur intégrée) ; système mono-tenant réservé à Caritas.

\section{Acteurs}

\begin{longtable}{@{}L{2.8cm}L{4.5cm}L{7.2cm}@{}}
\toprule
\textbf{Acteur} & \textbf{Objectif} & \textbf{Actions autorisées} \\
\midrule
\endhead
Administrateur & Qualité et bon fonctionnement des annonces & CRUD annonces, modération, gestion des comptes \\
Vendeur (bénévole) & Vendre rapidement les livres du stock & Ajouter, fixer un prix, modifier, supprimer ses annonces \\
Acheteur & Trouver et acheter un livre & Rechercher, filtrer, consulter, acheter via PostFinance \\
\bottomrule
\caption{Acteurs du système}
\end{longtable}

\noindent\textbf{Gouvernance :} la commanditaire valide les livrables de sprint ; Alexandre REY
conçoit et développe seul la solution ; l'encadrant pédagogique valide la conformité au
référentiel CDA.

\section{Spécifications fonctionnelles}

Chaque spécification (\texttt{SF-XX}) est reprise dans le plan de tests du mémoire pour
assurer la traçabilité besoin $\rightarrow$ test.

\begin{longtable}{@{}L{1.1cm}L{5.5cm}L{7.9cm}@{}}
\toprule
\textbf{Réf.} & \textbf{Spécification} & \textbf{Critère d'acceptation} \\
\midrule
\endhead
SF-01 & Authentification et rôles, mot de passe haché & Aucune route de gestion accessible sans authentification. \\
SF-02 & Ajout d'un livre par scan ISBN, saisie manuelle en secours & Livre reconnu ajouté en moins de 2 minutes. \\
SF-03 & Catalogue et recherche multicritère & Résultat cohérent en moins d'une seconde perçue. \\
SF-04 & Fiche détaillée & Informations essentielles visibles sans action supplémentaire (la liste de favoris a été reportée au backlog, hors MVP). \\
SF-05 & Panier avec réservation temporaire du stock & Livre jamais vendu deux fois ; panier abandonné libère le stock automatiquement. \\
SF-06 & Paiement CHF via PostFinance & Paiement réussi en sandbox déclenche le statut \og payée \fg{}. \\
SF-07 & Livraison restreinte à la Suisse & Seuls deux modes de livraison sont proposés (envoi postal, retrait en magasin), tous deux réservés au territoire suisse ; l'adresse de facturation n'est pas validée côté serveur (limite connue, voir recette). \\
SF-08 & Back-office bénévole & Gestion d'une annonce sans assistance technique. \\
SF-09 & Conformité RGPD / nLPD & Demande d'effacement supprime effectivement les données. \\
SF-10 & Accessibilité de base & Parcours critiques utilisables au clavier. \\
\bottomrule
\caption{Spécifications fonctionnelles et critères d'acceptation}
\end{longtable}

\section{Exigences non fonctionnelles et contraintes}

\begin{longtable}{@{}L{2.8cm}L{10.7cm}@{}}
\toprule
\textbf{Catégorie} & \textbf{Exigence} \\
\midrule
\endhead
Sécurité & Jeton signé à durée de vie limitée, mots de passe hachés, validation serveur, contrôle des permissions par rôle. \\
Devise / territoire & Transactions en CHF, livraison restreinte à la Suisse. \\
Disponibilité & Dégradation en saisie manuelle si la source de métadonnées externe est indisponible. \\
Performance & Catalogue et recherche réactifs (< 1 s perçue). \\
Ergonomie & Utilisable par des bénévoles non technophiles, sans formation. \\
Conformité légale & Respect du RGPD (UE) et de la nLPD (Suisse). \\
Maintenabilité & Architecture en couches (présentation, services, infrastructure). \\
Coût & Aucun abonnement SaaS récurrent, frais de paiement transparents. \\
\bottomrule
\caption{Exigences non fonctionnelles}
\end{longtable}

\noindent\textbf{Contraintes techniques :} MySQL exclusivement, backend FastAPI (Python) en
couches, frontend Next.js (React/TypeScript) 100\,\% typé, métadonnées via OpenLibrary,
déploiement continu par pipeline GitHub Actions (SSH, service systemd, sans conteneurisation),
tests automatisés avant déploiement.

\noindent\textbf{Contraintes organisationnelles :} développement solo en alternance
(environ 18h/sprint), projet de 3 mois en 6 sprints de 2 semaines, méthode Agile Scrum
simplifiée, aucun budget dédié.

\section{Livrables, planning et risques}

\noindent\textbf{Livrables :} code source (Git), application déployée, mémoire de projet,
présent cahier des charges, liste des compétences, support de soutenance, suite de tests
automatisés (CI), documentation API (Swagger/OpenAPI).

\noindent\textbf{Planning :} 6 sprints de 2 semaines — environnement \& modélisation (S1),
authentification (S2), ajout de livre / scan ISBN (S3), catalogue \& fiche livre (S4),
paiement (S5), back-office, tests et déploiement (S6).

\noindent\textbf{Principaux risques :} livres sans ISBN (mitigation : saisie manuelle) ;
vente simultanée du même exemplaire (mitigation : réservation temporaire de stock) ; dérive
de périmètre (mitigation : application stricte du MVP) ; indisponibilité de la source externe
de métadonnées (mitigation : saisie manuelle en secours).

\section{Recette}

La recette est réalisée par démonstration à la commanditaire, en fin de sprint puis en fin de
projet, sur la base des critères d'acceptation de la section 4.

\begin{longtable}{@{}L{1.3cm}L{8.5cm}L{3.9cm}@{}}
\toprule
\textbf{Réf.} & \textbf{Critère de réception} & \textbf{État} \\
\midrule
\endhead
SF-01 & Aucun accès non authentifié aux routes de gestion & Conforme \\
SF-02 & Ajout d'un livre reconnu en moins de 2 minutes & Conforme \\
SF-03 & Recherche multicritère fonctionnelle & Conforme \\
SF-04 & Fiche détaillée complète et lisible (favoris hors périmètre) & Conforme \\
SF-05 & Libération automatique du stock réservé & Conforme \\
SF-06 & Paiement CHF validé en environnement de test & Conforme avec réserve \\
SF-07 & Livraison restreinte à la Suisse (par choix du mode de livraison) & Conforme avec réserve \\
SF-08 & Gestion autonome des annonces par un bénévole & Conforme \\
SF-09 & Export / effacement des données personnelles & Conforme \\
SF-10 & Parcours critiques utilisables au clavier & Conforme avec réserve \\
\bottomrule
\caption{Grille de recette du MVP}
\end{longtable}

\noindent Le détail des tests et des réserves est documenté dans le chapitre \og Tests et
validation \fg{} du mémoire de projet.

\end{document}
